\chapter{Wishlists, Matching, Clubs \& Events}
\label{ch:wishlists-clubs}

\section{Wishlists: asking for a book that doesn't exist yet}

A wishlist entry (\pathname{wishlist/page.tsx}, Chapter~\ref{ch:pages}) is just a title, author, and genre a user wants but couldn't find on Browse. On its own it does nothing — it becomes useful only through the matching flow triggered from the \emph{other} side of the transaction.

\subsection*{The matching flow}

\begin{enumerate}
  \item User A adds \emph{"The Old Man and the Sea"} to their wishlist. Nothing happens yet — \code{matched\_book\_id} stays \code{NULL}.
  \item User B, independently, lists a copy of \emph{"The Old Man and The Sea"} (different capitalisation, doesn't matter) on \pathname{my-books}.
  \item \pathname{my-books/page.tsx} calls \code{match\_wishlists\_for\_book(title, owner\_id)} (Chapter~\ref{ch:functions}) right after the insert.
  \item The \code{pg\_trgm} similarity search finds User A's wishlist row (similarity comfortably above the 0.35 threshold despite the capitalisation difference), and the frontend updates that row's \code{matched\_book\_id} and fires a notification to User A.
\end{enumerate}

\begin{olknote}
Re-read Chapter~\ref{ch:functions}'s explanation of \emph{why} this needs to be a \code{SECURITY DEFINER} RPC rather than a plain query: the \code{wishlists} RLS policy (Chapter~\ref{ch:rls}) restricts a normal \code{SELECT} to \code{user\_id = auth.uid()}. User B's own session, running the match search, could never see User A's wishlist row through an ordinary query — the whole feature would be silently non-functional without the RPC's elevated privilege, which is exactly the bug this migration fixed.
\end{olknote}

The unique index on \code{(user\_id, lower(btrim(title)))} (Chapter~\ref{ch:schema}) exists to stop a second, independent problem: without it, adding the same title twice (perhaps after a page refresh double-submits a form) would previously have produced two wishlist rows, and a later match would notify the user twice for the same book.

\section{Clubs: geography plus a trust gate}

Clubs are the only feature in the schema besides books themselves that carries its own latitude/longitude (Chapter~\ref{ch:schema}) and its own nearby-search RPC, \code{get\_clubs\_nearby} (Chapter~\ref{ch:functions}) — the same haversine pattern, applied to a different table, ordered by distance then by \code{member\_count}.

\subsection*{Creating a club is submit-then-review, not instant self-service}

\pathname{clubs/create/page.tsx} (Chapter~\ref{ch:pages}) does not insert into \code{clubs} directly. It inserts into \code{club\_requests} — the name, description, interests, and the rest of a club's future shape, plus a \code{status} that starts at \code{pending} — and a real \code{clubs} row only ever comes into existence once a moderator approves that request.

\begin{olknote}
This wasn't the original design. Clubs used to be a plain self-service \code{INSERT} into \code{clubs}, gated only by an eligibility trigger on that table. \pathname{20260718150000\_club\_requests\_and\_review\_workflow.sql} replaced that with the submit-then-review flow described here, and — deliberately — \code{DROP}ped the \code{clubs} table's INSERT policy for regular users entirely (Chapter~\ref{ch:rls}). \code{approve\_club\_request} below is now the \emph{only} path that ever inserts a \code{clubs} row for a normal user; it runs \code{SECURITY DEFINER}, so it doesn't need a policy of its own to do it.
\end{olknote}

\subsection*{The eligibility bar, and one function that owns it}

Submitting a request requires 5+ completed exchanges (as either the book's owner or the requester — Chapter~\ref{ch:schema}) and zero reports against the requester. That rule used to be implemented three separate times: once in a \code{BEFORE INSERT} trigger on \code{clubs}, again nearly verbatim in a second trigger added on \code{club\_requests} when the review workflow shipped, and a third time as two client-side count queries in \pathname{clubs/create/page.tsx} just to paint the eligibility checklist. All three could drift out of sync — the two trigger copies even used different column names (\code{creator\_id} vs. \code{requester\_id}) for the same underlying check — which meant a threshold change had to be applied by hand in three unrelated files, and it was easy to update one copy and forget the others.

\pathname{20260802055510\_refactor\_club\_eligibility\_shared\_function.sql} consolidated all three into a single function, \code{get\_club\_eligibility(user\_id)} (Chapter~\ref{ch:functions}), returning \code{(eligible, completed\_exchanges, report\_count)} for whichever user ID it's given. Both triggers now call it instead of re-deriving the counts themselves: \code{check\_club\_request\_eligibility} on \code{club\_requests} is the one that actually gates a normal user today, while \code{check\_club\_creation\_eligibility} on \code{clubs} is only reached via \code{approve\_club\_request}'s own insert — a redundant second check at approval time, kept as defence in depth rather than deleted. The client reads the same numbers through a second, narrower RPC, \code{my\_club\_eligibility()}, which wraps \code{get\_club\_eligibility(auth.uid())}.

\begin{olknote}
\code{get\_club\_eligibility} itself is \emph{not} granted to \code{authenticated} (\code{REVOKE ALL ... FROM public}) — it takes an arbitrary \code{user\_id}, and its \code{report\_count} would let any logged-in caller check whether some \emph{other} user has been reported, simply by guessing their UUID. Only \code{my\_club\_eligibility()} — which hard-codes the caller to \code{auth.uid()} — is exposed to PostgREST. The two triggers call \code{get\_club\_eligibility} directly, which is fine: they already run \code{SECURITY DEFINER}, so they aren't subject to the grant at all.
\end{olknote}

\begin{olkhowto}
Raising the bar to, say, 10 exchanges is now a one-function edit. Run \code{npx supabase migration new raise\_club\_threshold}, and in the new migration file put a \code{CREATE OR REPLACE FUNCTION public.get\_club\_eligibility(...)} whose final \code{RETURN QUERY SELECT} checks \code{v\_completed\_exchanges >= 10} instead of \code{>= 5}. Never hand-edit \pathname{supabase/schema.sql} directly — it's a dump, not a source (Chapter~\ref{ch:local-env}). Nothing else needs to change: both triggers and the client-facing checklist all read the new threshold the moment the migration lands, since none of them hard-code the number anymore.
\end{olkhowto}

\subsection*{Review: a moderator's call, always logged}

\pathname{admin/club-requests/page.tsx} (Chapter~\ref{ch:pages}) lists pending requests for any \code{moderator} or above. Approving calls \code{approve\_club\_request(request\_id, note)} (Chapter~\ref{ch:functions}), which — inside one \code{SECURITY DEFINER} function, so it can't partially apply — inserts the real \code{clubs} row, adds the requester as its first \code{club\_members} row, and marks the request \code{approved} with the reviewer and their note. Rejecting just flips \code{status} to \code{rejected}; no \code{clubs} row is ever created. Both actions run through the same \code{withAdminAction('moderator', ...)} wrapper as every other admin action (Chapter~\ref{ch:admin}) — guarded, written to \code{admin\_audit\_log}, and followed by a notification to the requester either way.

A requester can withdraw their own request while it's still \code{pending} — a \code{DELETE} policy scoped to \code{auth.uid() = requester\_id AND status = 'pending'} — the only regular-user write \code{club\_requests} allows directly. There is deliberately no \code{UPDATE} policy on the table at all: even a moderator's own session cannot flip \code{status} by hand, only \code{approve\_club\_request} and \code{reject\_club\_request} can, so every approval or rejection is guaranteed to go through the audited path.

\begin{olktask}
Submit a club request from a test account with fewer than 5 completed exchanges, confirming the UI's checklist correctly disables the form — then try inserting into \code{club\_requests} directly via \code{POST /rest/v1/club\_requests} with that account's own access token, and confirm it still comes back rejected by \code{check\_club\_request\_eligibility}. Then, from an eligible account, submit a request and approve it from a \code{moderator} test account; confirm a real \code{clubs} row and a \code{club\_members} row both exist afterward, and that the requester received a notification. Same "is it a real gate or just a hidden button" habit this chapter asks for club creation generally, now applied end to end across the request-and-review flow.
\end{olktask}

\subsection*{Membership and posts}

Joining is a simple \code{club\_members} insert, kept honest by \code{trg\_club\_member\_count} (Chapter~\ref{ch:functions}) rather than client-side arithmetic. \code{removeClubMember} in admin actions (Chapter~\ref{ch:admin}) used to redundantly re-count on top of the trigger, introducing an off-by-one on every admin-initiated removal; that manual re-count has been removed.

Posting to a club is creator-only (Chapter~\ref{ch:rls}); a new post fans out one notification per member via a single batched \code{.insert(notifications[])} call (Chapter~\ref{ch:pages}) — the same N+1-avoidance discipline seen in the messages list (Chapter~\ref{ch:pages}) and the admin broadcast action (Chapter~\ref{ch:admin}).

\subsection*{Deactivation is soft, always}

Neither a club's creator nor an admin can hard-delete a club through the app layer — \code{active = false} is the only operation exposed (\pathname{clubs/[id]/page.tsx}'s "delete," and the admin \code{deactivateClub}/\code{reactivateClub} actions, Chapter~\ref{ch:admin}). Only a \code{super\_admin}, and only directly via SQL, could actually remove a \code{clubs} row (Chapter~\ref{ch:rls}'s DELETE policy) — there is no UI path to a hard delete at all, by design.

\begin{olktask}
Add a wishlist entry for a title that doesn't exist yet, then list a book with a \emph{deliberately misspelled} version of that title from a second test account (swap two letters, or use a different capitalisation) and confirm the match still fires. Then try a title that only shares one common word (e.g. "The Old Man" vs. "Old Man River") and confirm it does \emph{not} match — this is the fastest way to build a feel for where the 0.35 similarity threshold actually sits in practice, beyond just reading the number.
\end{olktask}

\section{Events: club-scoped, but publicly discoverable}

A club can schedule events (\code{club\_events}, plus one \code{event\_rsvps} row per attendee — Chapter~\ref{ch:schema}): a title, description, start/end time, either an in-person location or an online meeting link, an optional attendee cap, and a per-event \code{visibility} of \code{public} or \code{members\_only}.

\begin{olknote}
\textbf{\code{visibility} gates the RSVP action, not discovery.} Every event stays listed and viewable by anyone — on the club's own page, on the cross-club \pathname{/events} directory, and in the landing page's "Upcoming Events" preview — regardless of the viewer's club membership, mirroring how clubs themselves are already publicly browsable (Chapter~\ref{ch:rls}). A \code{members\_only} event just shows "Join the club to RSVP" instead of an RSVP button for a non-member. Two different features, deliberately not conflated: hiding an event outright would need a second, separate flag.
\end{olknote}

Creating an event is gated exactly like posting an announcement (Chapter~\ref{ch:rls}) — only the club's creator — rather than the club-creation eligibility rule earlier in this chapter; a club's creator has already cleared that bar once, so there is no second trigger to re-derive it for events. \code{trg\_event\_attendee\_count} keeps \code{club\_events.attendee\_count} in sync with real \code{event\_rsvps} rows, the same atomic-trigger shape as \code{trg\_club\_member\_count} above (Chapter~\ref{ch:functions}).

\begin{olkgotcha}
The RSVP \code{INSERT} policy's \code{visibility} check was added alongside a second, independent condition — \code{capacity IS NULL OR attendee\_count < capacity} — but that second condition didn't ship on day one. The "Event full" state was originally enforced only by disabling the RSVP button once \code{attendee\_count >= capacity} in \pathname{events/[id]/page.tsx}; a direct \code{POST /rest/v1/event\_rsvps} call could still RSVP past a full event, confirmed by testing exactly that against a capacity-1 event before the policy was corrected. Same lesson as the club-creation gate earlier in this chapter, one feature later: a disabled button is a UX nicety, not a security boundary — only the policy is.
\end{olkgotcha}

Event notifications reuse the club-announcement broadcast shape: creating an event fans out an \code{event\_created} notification to every current member (Chapter~\ref{ch:server-actions}), routed through the same \code{can\_notify} guard with a new branch verifying the caller is the event's own creator and the target is a real club member. There is deliberately no per-RSVP notification to the organiser — at any real scale that would be pure noise, so the attendee list on the event and admin pages covers that need instead.

The cross-club directory (\pathname{/events}) uses \code{get\_events\_nearby}, the same haversine-and-\code{NULLS LAST} pattern as \code{get\_clubs\_nearby} (Chapter~\ref{ch:functions}), falling back to a plain \code{created\_at}-ordered query when the viewer has no location set. A zero-dependency Route Handler, \pathname{api/events/[id]/ics/route.ts}, generates an RFC5545 calendar file on demand for "Add to calendar" — no external library, just a template string with the handful of characters (\code{,}, \code{;}, \code{\textbackslash}) that need escaping inside an ICS field.

\begin{olktask}
Create a club event with \code{capacity} set to 1, RSVP to it yourself, then — from a second test account that is \emph{not} a member of that club — try to RSVP through the UI (should be blocked by the disabled button) and then directly via a \code{POST} to \pathname{/rest/v1/event\_rsvps} with that account's own access token (should come back \code{403}). This is the same "is it a real gate or just a hidden button" habit Chapter~\ref{ch:rls}'s task asks you to build for club creation, applied to the newest feature in the schema.
\end{olktask}
